\documentclass[11pt]{article}
\usepackage{icmlko}
\begin{document}

\title{축이 많은 행은 원래 쉬운 행이었다 --- $+0.10$ 중 축의 몫은
$5$분의 $1$}
\author{월드모델 실험실}
\date{노트 480 · 2026년 8월 3일}
\maketitle

\begin{abstract}
노트 479 가 구조적 사실 하나를 냈다 --- \textbf{값 열 $37$개 중 덮음
$95\%$ 를 넘는 것은 셋뿐}이고 나머지 $34$개는 행의 절반 이상에서 없다.
그러면 \emph{없는 축은 얼마나 비싼가}를 물어야 한다.

도메인 \emph{안}에서 유보 행을 관측 축 수로 반 갈랐다. 큰 수가 나왔다.

\begin{center}
\begin{tabular}{lrr}
\hline
 & 축 많은 절반 & 축 적은 절반 \\
\hline
가중평균 & $\mathbf{0.6161}$ & $\mathbf{0.5133}$ \\
\hline
차 & \multicolumn{2}{c}{$\mathbf{+0.1028}$ ($4/5$ 도메인)} \\
\hline
\end{tabular}
\end{center}

만화 $+0.2640$ · 모바일 $+0.2244$ --- 이 사이클의 어떤 채택보다 한
자릿수 크다. \textbf{여기서 멈추면 ``태거를 늘리면 크게 번다''가 된다.}

\textbf{안 멈췄다.} 축이 많은 행은 자료 출처가 많이 붙은 행이고, 그런
작품은 대개 유명하고 라벨이 덜 시끄럽다. 그러면 \emph{축이 도운 게
아니라 행이 원래 쉬운 것}이다. 가르는 법은 하나뿐이다 --- \textbf{덮음
$95\%$ 넘는 축 셋만 쓰는 예보로 같은 두 절반을 채점한다.} 축 셋은 두
절반에 (거의) 똑같이 있으므로, 그래도 격차가 남으면 행 탓이다.

\begin{center}
\begin{tabular}{lr}
\hline
격차 & \\
\hline
축 전부 & $+0.1028$ \\
\textbf{축 셋만} & $\mathbf{+0.0808}$ \\
\hline
줄어든 비율 & $\mathbf{21\%}$ \\
\hline
\end{tabular}
\end{center}

\textbf{$79\%$ 는 행이 원래 쉬운 것이다.}
\end{abstract}

\begin{figure}[t]
\centering
\includegraphics[width=\textwidth]{figs/f1.pdf}
\caption{\textbf{가} --- 축을 셋으로 줄여도 격차가 거의 그대로다.
\textbf{나} --- $+0.1028$ 의 분해.}
\end{figure}

\section{그래서 태거의 기대값이 내려간다}

노트 $444 \cdot 454 \cdot 471$ 이 반복해서 \emph{태거 갈래}를 가리켰다
--- CN 의 낮은 천장도, 시장팝업의 열여덟 칸도, 펀딩의 조회표도 결국
``축을 더 만들어라''였다. 이 노트는 그 갈래의 값을 \textbf{내린다}.

행에 축을 하나 더 붙여서 얻는 것은 관찰된 격차의 $21\%$ 다. 절반을
``적은 쪽''에서 ``많은 쪽''으로 옮긴다 해도 $+0.1028 \times 0.21 \approx
+0.022$ 이고, 그것도 \emph{그 절반에 대해서만}이다. \textbf{여전히
채택 하나($+0.004$)보다 크지만 $+0.10$ 은 아니다.}

\section{정직하게 적을 것 셋}

\textbf{하나 --- 대조가 완벽하지 않다.} 덮음 $95\%$ 는 \emph{전체}
기준이고 도메인마다 다르다. 모바일에서 축 셋만 쓴 격차가 $+0.4612$ 로
\emph{전체 축일 때($+0.2244$)보다 크다} --- 그 도메인에서는 셋 중 하나가
두 절반에 고르게 없다는 뜻이다. 대조가 그만큼 샜다.

\textbf{둘 --- 도메인마다 크게 다르다.} 애니는 격차가 $-0.1004$ 로
\emph{반대}이고 축 셋만 쓰면 $-0.2257$ 로 더 벌어진다. 가중평균이 그것을
덮고 있다.

\textbf{셋 --- 축 수는 라벨과도 조금 붙는다}(가중평균 $+0.088$ · 게임은
$+0.39$). 두 절반의 라벨 분포가 완전히 같지는 않다.

\section*{조항 --- ``더 있으면 더 잘한다''는 대조가 필요하다}

\begin{quote}
\textbf{조항} --- 어떤 자료가 있는 행과 없는 행을 갈라 성능을 견줄 때는,
\textbf{그 자료를 안 쓰는 예보로 같은 두 무리를 채점해} 행 자체의
난이도를 먼저 뺀다. 안 빼면 \emph{자료의 값과 행의 값을 더해서 자료에
붙인다.}
\end{quote}

\section*{남은 것}

\begin{itemize}
\item \textbf{애니는 왜 반대인가} --- 축이 적은 절반이 더 잘 맞는
유일한 도메인이고, 축 셋만 쓰면 더 벌어진다.
\item \textbf{``원래 쉬운 행''의 정체} --- 유명해서인가, 라벨이 덜
시끄러워서인가, 아니면 셋 중 하나가 그 행들에서 더 잘 듣는 것인가.
안 갈랐다.
\item \textbf{덮음 $95\%$ 를 도메인마다 다시} --- 이번 대조가 샌 이유다.
\end{itemize}

\end{document}
